\clearpage
\VerdierPhysicalPageAnchor{0194}%
\begingroup\small\setlength{\parskip}{0pt}\setlength{\emergencystretch}{1.5em}
\setlength{\abovedisplayskip}{4pt plus 2pt minus 2pt}%
\setlength{\belowdisplayskip}{4pt plus 2pt minus 2pt}%

\VerdierSemanticAnchor{ch3d-2-2-6-proof-cone-cont}
（$L_0$ は $c(m)$ の次数 $0$ の成分である。）$c(m)$ のコホモロジーは
$-1$ 以外のすべての次数で零であり、次数 $-1$ で $X$ に等しい。
射 $L_n\longrightarrow L_n\oplus Z_{n-1}$ の行列は
$\begin{pmatrix}\operatorname{id}\\ *\end{pmatrix}$ の形をする。従って、この射はレトラクションをもつ。
したがって、射 $v\colon Z_{n-1}\longrightarrow L_{n-1}\oplus L'_{n-2}$ であって、複体

\VerdierSemanticAnchor{ch3d-2-2-6-reduced-cone}
\[
\begin{aligned}
\cdots\longrightarrow0\longrightarrow0\longrightarrow Z_{n-1}
\xrightarrow{v}L_{n-1}\oplus L'_{n-2}\longrightarrow\cdots
\longrightarrow L_1\oplus L'_0\\
\longrightarrow L_0\longrightarrow0\longrightarrow\cdots
\end{aligned}
\]
が $c(m)$ と擬同型になるものが存在する。一方、性質
$(E_4)$ により、$L_1\oplus L'_0\longrightarrow L_0$ の核は $M$ に属する。
この核を $L''_1$ と記す。複体

\VerdierSemanticAnchor{ch3d-2-2-6-resolution-x}
\[
\begin{aligned}
0\longrightarrow0\longrightarrow Z_{n-1}
\longrightarrow L_{n-1}\oplus L'_{n-2}\longrightarrow\cdots
\longrightarrow L_2\oplus L'_1\\
\longrightarrow L''_1\longrightarrow X\longrightarrow0\longrightarrow\cdots
\end{aligned}
\]
は非輪状である。こうして $X$ の長さ $n-1$ の分解を構成したので、
$M\text{-}\dim(X)\leq n-1$ である。一方、明らかに
$M\text{-}\dim(X)\geq M\text{-}\dim(Y)-1$ である。

\VerdierSemanticAnchor{ch3d-2-2-6-conclusion}
これで補題 (2.2.6) と命題 (2.2.7) の証明が完了した。

\VerdierSemanticAnchor{ch3d-2-3}
\medskip\noindent\textbf{2.3. 分解に関するいくつかの補題。}

\VerdierSemanticAnchor{ch3d-2-3-1}
\medskip\noindent\textbf{命題 2.3.1.} \textbf{[22]}\quad
\textit{a) $M$ を $\mathcal A$ の対象の集合で、性質 $(E_1)$, $(E_2)$, $(E_3)$
（2.2.1）を満たすものとする。$Y\in\operatorname{Ob}\operatorname{comp}^{-}(\mathcal A)$
を、その各成分が $M^\wedge$（2.2.1）に属する複体とし、
$m\colon X\longrightarrow Y$ を擬同型とする。このとき、複体
$L^\bullet\in\operatorname{Ob}\operatorname{comp}^{-}(\mathcal A)$ で各成分が $M$ に属するものと、擬同型
$w\colon L^\bullet\longrightarrow X$ が存在する。}

\noindent\textit{b) 逆に、$M$ を $\mathcal A$ の対象の集合で、性質 $(E_1)$ を満たすものとする。
任意の複体 $Y\in\operatorname{Ob}\operatorname{comp}^{-}(\mathcal A)$ で各成分が $M$ に属するものと任意の擬同型
$m\colon X\longrightarrow Y$ に対して、複体 $L^\bullet$ で各成分が $M$ に属するものと擬同型
$w\colon L^\bullet\longrightarrow X$ が存在すると仮定する。このとき $M$ は性質
$(E_2)$ と $(E_3)$ を満たす。}

\VerdierSemanticAnchor{ch3d-2-3-1-proof}
a) を証明する。まず、複体
$X'\in\operatorname{Ob}\operatorname{comp}^{-}(\mathcal A)$ と擬同型
$X'\longrightarrow X$ が存在する（1.2.1）。したがって、$X$ が
$\operatorname{comp}^{-}(\mathcal A)$ の対象である場合に帰着できる。
(2.1.4) と (2.1.5) により、射の列

\VerdierSemanticAnchor{ch3d-2-3-1-xn-sequence}
\[
u_n\colon X_n\longrightarrow X_{n-1},
\qquad n\in\mathbf Z,
\]
が存在する。
\endgroup
